\documentclass[10pt,a4paper]{amsart}
\usepackage[margin=19mm]{geometry}
\usepackage{amsmath,amssymb,mathtools}
\usepackage[scr=rsfs,cal=euler]{mathalfa}
\usepackage{xfrac}
\usepackage[unicode,hidelinks,bookmarks=false]{hyperref}
\newcommand{\ie}{i.e.\ }
\newcommand{\cf}{cf.\ }
\newcommand{\resp}{respectively\ }
\newcommand{\Subsection}{\subsection}
\providecommand{\nobreakdash}{\nobreak\hbox{-}}
\setlength{\emergencystretch}{2em}
\multlinegap0pt
\usepackage{bm}
\usepackage{bbm}
\usepackage{dsfont}
\usepackage{xspace}
\usepackage{cancel}
\usepackage{mathtools}
\usepackage{amsthm}
\makeatletter
\@ifundefined{proposition}{\newtheorem{proposition}{Proposition}}{}
\@ifundefined{remark}{\newtheorem{remark}{Remark}}{}
\makeatother
\makeatletter
\newcommand{\R}{\mathbb R}
\expandafter\ifx\csname rcases\endcsname\relax
\newenvironment{rcases}
%   {\left.\begin{aligned}}
%   {\end{aligned}\right\rbrace}
%%
\fi
\makeatother
\title[A convex variational principle for the necessary conditions ]{A convex variational principle for the necessary conditions of classical optimal control}
\author{Amit Acharya, Janusz Ginster}
\date{Source: arXiv:2502.15973v2 (CC BY 4.0)}
\begin{document}
\maketitle

\noindent\emph{Source and licence.} A convex variational principle for the necessary conditions of classical optimal control. Version arXiv:2502.15973v2 (CC BY 4.0), distributed under the Creative Commons Attribution 4.0 International licence, as recorded in the arXiv licence metadata for this exact version and shown on the abstract page https://arxiv.org/abs/2502.15973v2. Full rights notice: licenses/arxiv-2502.15973-CC-BY-4.0.txt.\par\smallskip

\noindent\textit{Editorial scope.} Counted unit: the Remark whose source label is \texttt{None} in the article (source0.tex lines 647--672, source label \texttt{None}), delivered together with the article's own complete original proof of it (source0.tex lines 675--797, one \texttt{proof} environment of 123 lines). No mathematics is written, repaired or re-derived: statement and proof are the article's own text line for line. Required context retained uncounted: prop: lb g (lines 605--614); eq: defR (lines 632--634); eq: invert (lines 636--638). Every definition, display and label of those lines is retained unchanged. Omitted from this package and not redistributed: 1--604, 615--631, 635--635, 639--646, 673--674, 798--1253. Nothing inside a retained excerpt is omitted or altered: every retained body is delivered exactly as the article's lines are written, so no omission receipt of any kind accompanies this package. Citations: the retained excerpts carry no citation command at all, so no bibliography entry of the article is transcribed and no reference is redistributed. Reference adaptation: none was needed and none was made; every reference inside the delivered unit resolves inside it, so no unresolved reference or citation is produced. Printed slips kept literal: no printed word, symbol, hypothesis or number of the counted statement or of its proof was altered. Layout and class: the article's own class is \texttt{article} with packages including graphicx, fancyhdr, isomath, mathtools, amsbsy, amsmath, amsthm, amssymb, amscd, amsfonts, bigints, verbatim, euscript, alltt; the wrapper is the lane's pinned \texttt{amsart} editorial wrapper at 10pt on A4, which restates the theorem-like environments the retained text needs and adds the article's own macro definitions that the retained text needs (\texttt{\textbackslash R}), with extra wrapper packages (bm, bbm, dsfont, xspace, cancel, mathtools). The delivered numbering is the unit's own. Assets: no retained span contains a figure environment, a graphics-inclusion command, a PDF-page inclusion, a table or a TikZ picture; no image file is copied into this package, the package declares no asset dependency and no image file is redistributed with it. Licence: version arXiv:2502.15973v2 of this article is distributed under the Creative Commons Attribution 4.0 International licence, as recorded in the arXiv licence metadata for this exact version and shown on the abstract page https://arxiv.org/abs/2502.15973v2. A full rights notice is delivered at licenses/arxiv-2502.15973-CC-BY-4.0.txt. Characters: every retained excerpt is pure ASCII, so the delivered engine, XeLaTeX with its default Latin Modern text font, reports no missing character.
\begin{proposition}\label{prop: lb g}
    It holds for all $(\gamma,\alpha,\lambda,\beta,\mu) \in \R^n \times \R^n \times \R^n \times \R^n \times \R^m$ that
    \begin{align*}
    g(\gamma,\alpha,\lambda,\beta,\mu) \geq &\frac{\left| \alpha + M^T \gamma + B^T \lambda \right|^2}{2 (a_x + |F| |\gamma| + |F| |\lambda|)}  + \frac{\left| \beta - M \lambda + N \mu \right|^2}{2 ( a_p + |F| |\lambda|)}  + \frac{\left| N^T \gamma + C \mu \right|^2}{2a_u} .
    \end{align*}
    For $F = 0$ it holds 
    \[
        g(\gamma,\alpha,\lambda,\beta,\mu) = \frac{\left| \alpha + M^T \gamma + B^T \lambda \right|^2}{2 a_x}  + \frac{\left| \beta - M \lambda + N \mu \right|^2}{2  a_p }  + \frac{\left| N^T \gamma + C \mu \right|^2}{2a_u} .
    \]
\end{proposition}

\begin{equation}\label{eq: defR}
R := \begin{pmatrix} -M^T && - B^T \\ N C^{-1}N^T &&  M \end{pmatrix} \in \R^{2n \times 2n}.
\end{equation}

\begin{equation} \label{eq: invert}
\R^n \ni x \rightarrow \pi_1 (e^{TR}  \iota(x)) + G \pi_2 (e^{TR} \iota(x)) 
\end{equation}

\begin{remark}
    Let us briefly comment on the invertibility condition \eqref{eq: invert}. We write 
    \[
    \rho := \dot{\gamma} + M^T \gamma + B^T \lambda \text{ and } \sigma := \dot{\lambda} - M \lambda + N\mu = \dot{\lambda} - M \lambda - N C^{-1} N^T\gamma + NC^{-1} ( N^T \gamma + C \mu).
    \]
    Hence, the functions $\gamma, \lambda$ solve the ODE
    \[
    \begin{pmatrix}
        \dot{\gamma} \\ \dot{\lambda}
    \end{pmatrix} = R \begin{pmatrix}
        \gamma \\ \lambda
    \end{pmatrix} + \begin{pmatrix}
        \rho \\ \sigma - NC^{-1}(N^T \gamma + C \mu)
    \end{pmatrix},
    \]
    where $R$ is the matrix in \eqref{eq: defR}. 
    By Proposition \ref{prop: lb g} it seems reasonable that for admissible functions $(\gamma,\lambda,\mu) \in \mathcal{A}$ with a bounded energy $\tilde{S}_Q[\gamma,\lambda,\mu]$ one might hope to control the functions $\rho$, $\sigma$ and $N^T \gamma + C \mu$ so that it holds approximately
    \[
    \begin{pmatrix}
        \gamma(t) \\ \lambda(t)
    \end{pmatrix} \approx e^{tR} \begin{pmatrix}
        \gamma(0) \\ \lambda(0)
    \end{pmatrix}.
    \]
    Then the invertibility of the mapping in \eqref{eq: invert} is exactly the condition that allows to control $\gamma(0)$ through $\lambda(0) = \lambda_0$ using that $\gamma(T) = G \lambda(T)$. In turn this will imply bounds on the functions $\gamma$ and $\lambda$ in $H^1((0,T);\R^n)$.
\end{remark}

\begin{proof}
Let $(\gamma^{(k)},\lambda^{(k)}, \mu^{(k)})_k \subseteq \mathcal{A}$ be a sequence satisfying $\sup_k \tilde{S}_Q[\gamma^{(k)},\lambda^{(k)}, \mu^{(k)}] \leq K$, for some $K>0$.
By the usual weak compactness properties of the spaces $H^1$ and $L^2$ it suffices to prove that the sequence $(\gamma^{(k)},\lambda^{(k)}, \mu^{(k)})_k$ lies in a bounded subset of $H^1(0,T; \R^m) \times H^1(0,T; \R^m) \times L^2((0,T); \R^m)$.

Throughout the proof $\overline{c}>0$ will denote a constant which does not depend on $\gamma^{(k)}$, $\lambda^{(k)}$ nor $\mu^{(k)}$ but may change from line to line.

Defining $\sigma^{(k)} = \dot{\gamma}^{(k)} + M^T\gamma^{(k)} + B^T \lambda^{(k)}$ and $\rho^{(k)} = \dot{\lambda}^{(k)} - M \lambda^{(k)}  + N \mu^{(k)}$ we obtain from Proposition \ref{prop: lb g} for a constant $c>0$ (depending on $a_x, a_p$ and $|F|$)
\begin{align}
&K \geq \tilde{S}_Q[\gamma^{(k)},\lambda^{(k)}, \mu^{(k)}]  \label{eq: est coer dual} \\
\geq &\int_0^T \left( \frac{|\sigma^{(k)}|^2}{c(1 + |\lambda^{(k)}| + |\gamma^{(k)}|)} + \frac{|\rho^{(k)} |^2}{c(1 + |\lambda^{(k)}| + |\gamma^{(k)}|)} + \frac1{2a_u} |N^T\gamma^{(k)} + C \mu^{(k)}|^2 + A\cdot \gamma^{(k)} \right)  \, dt - \gamma^{(k)}(0)\cdot x^0 \nonumber \\
\geq &\int_0^T \left( \frac{|\sigma^{(k)}|^2}{c(1 + |\lambda^{(k)}| +  |\gamma^{(k)}|)} + \frac{|\rho^{(k)} |^2}{c(1 + |\lambda^{(k)}| + |\gamma^{(k)}|)} + \frac1{2a_u} |N^T\gamma^{(k)} + C \mu^{(k)}|^2 \right)  \, dt \nonumber  \\
& - \|A\|_{L^1} \| \gamma^{(k)}\|_{L^{\infty}} - |\gamma^{(k)}(0)| \, |x^0| \nonumber \\
\geq &\int_0^T \left( \frac{|\sigma^{(k)}|^2}{c(1 + |\lambda^{(k)}| +  |\gamma^{(k)}|)} + \frac{|\rho^{(k)} |^2}{c(1 + |\lambda^{(k)}| +  |\gamma^{(k)}|)} + \frac1{2a_u} |N^T\gamma^{(k)} + C \mu^{(k)}|^2 \right)  \, dt - \delta \| \gamma^{(k)} \|_{L^{\infty}}. \nonumber
\end{align}
Next, we write
\begin{equation}\label{eq: rep derivative}
\begin{pmatrix}
    \dot{\gamma}^{(k)} \\ \dot{\lambda}^{(k)}
\end{pmatrix} = \underbrace{\begin{pmatrix} -M^T && - B^T \\ N C^{-1}N^T &&  M \end{pmatrix}}_{= R} \begin{pmatrix} \gamma^{(k)} \\ \lambda^{(k)} \end{pmatrix} - \begin{pmatrix} 0 \\ N C^{-1} ( N^T\gamma^{(k)} + C\mu^{(k)}) \end{pmatrix} +  \begin{pmatrix} \sigma^{(k)} \\ \rho^{(k)} \end{pmatrix}.
\end{equation}
This ODE is solved uniquely by 

\begin{equation} 
\begin{pmatrix} \gamma^{(k)}(t) \\ \lambda^{(k)}(t) \end{pmatrix} = e^{tR} \begin{pmatrix} \gamma^{(k)}(0) \\ \lambda^{(k)}(0) \end{pmatrix} + e^{tR} \int_0^t e^{-sR} \begin{pmatrix} \sigma^{(k)}(s) \\ - N^TC^{-1} ( N \gamma^{(k)}(s) + C\mu^{(k)}(s)) + \rho^{(k)}(s) \end{pmatrix} \, ds \label{eq: sol ode}
%&= e^{(t-T)R}\begin{pmatrix} \gamma^{(k)}(T) \\ \lambda^{(k)}(T) \end{pmatrix} - e^{tR} \int_t^T e^{-sR} \begin{pmatrix} \sigma^{(k)}(s) \\ -N C^{-1} ( N^T\gamma^{(k)}(s) + C\mu^{(k)}(s)) + \rho^{(k)}(s) \end{pmatrix} \, ds.\label{eq: sol ode 2}
\end{equation}
Writing the initial condition in the form $\begin{pmatrix} \gamma^{(k)}(0)& \lambda^{(k)}(0)\end{pmatrix}^T$ = $\begin{pmatrix} 0 & \lambda^{(k)}(0)\end{pmatrix}^T + \begin{pmatrix} \gamma^{(k)}(0) & 0\end{pmatrix}^T$ and using $\gamma^{(k)}(T) = G \lambda^{(k)}(T)$ and $\lambda^{(k)}(0) = \lambda^0$ in \eqref{eq: sol ode} with some rearrangement of terms, we find that
\begin{equation}
\begin{aligned}
&    \pi_1 (e^{TR} \iota (\gamma^{(k)}(0))) - G \pi_2 (e^{TR} \iota(\gamma^{(k)}(0))) \\ 
= & - (\pi_1 - G \circ \pi_2) \left( e^{TR} \begin{pmatrix} 0 \\ \lambda^0 \end{pmatrix} + e^{TR} \int_0^T e^{-sR} \begin{pmatrix} \sigma^{(k)}(s) \\ - N C^{-1} ( N^T\gamma^{(k)}(s) + C\mu^{(k)}(s)) + \rho^{(k)}(s) \end{pmatrix} \, ds \right).  
\end{aligned}
\end{equation}
By the invertibility of the mapping in \eqref{eq: invert} we obtain 
\begin{align}\label{eq: est gamma0}
    |\gamma^{(k)}(0)| 
    \leq \overline{c} \left( |\lambda^0| + \int_0^T |\sigma^{(k)}| + |\rho^{(k)}| + |N^T\gamma^{(k)} + C \mu^{(k)}| \, ds \right)
\end{align}
Combining this with \eqref{eq: sol ode}, we find using \eqref{eq: est coer dual} and H\"older's inequality 
\begin{align*}
    &\| \gamma^{(k)} \|_{L^{\infty}} + \| \lambda^{(k)} \|_{L^{\infty}} \\
    \leq &\overline{c} \left( |\gamma^{(k)}(0)| +  |\lambda^0| + \int_0^T |\sigma^{(k)}| + |\rho^{(k)}| + |N^T\gamma^{(k)} + C \mu^{(k)}| \, ds \right) \\
    \leq &\overline{c} \left(|\lambda^0| + \| N^T \gamma^{(k)} + C \mu^{(k)}\|_{L^2} +  \left( \int_0^T \frac{|\sigma^{(k)}|^2 + |\rho^{(k)}|^2 }{1+|\lambda^{(k)}| + |\gamma^{(k)}|} \, dt \right)^{1/2} (1 + \| \lambda^{(k)} \|_{L^{\infty}} + \| \gamma^{(k)} \|_{L^{\infty}})^{1/2} \right)
    \end{align*}
Using \eqref{eq: est coer dual} and Young's inequality it follows
    \begin{align*}
        &\| \gamma^{(k)} \|_{L^{\infty}} + \| \lambda^{(k)} \|_{L^{\infty}} \\
    \leq \  &\overline{c} \left(|\lambda^0| + (K + \delta \| \gamma^{(k)}\|_{L^{\infty}})^{1/2} +  \left( K + \delta \| \gamma^{(k)}\|_{L^{\infty}} \right)^{1/2}  (1 + \| \lambda^{(k)} \|_{L^{\infty}} + \| \gamma^{(k)} \|_{L^{\infty}})^{1/2}\right) \\
\leq \  &\overline{c} \left( |\lambda^0| + 1 +  K + \delta \| \gamma^{(k)}\|_{L^{\infty}} + \delta^{-1/2}( K + \delta \| \gamma^{(k)} \|_{L^{\infty}} ) + \delta^{1/2} ( 1 + \| \lambda^{(k)} \|_{L^{\infty}} + \| \gamma^{(k)} \|_{L^{\infty}}) \right).
\end{align*}
Hence, if 
\[
\overline{c}(2\delta^{1/2}+\delta) \leq \frac12,
\]
we find 
\[
        \| \gamma^{(k)} \|_{L^{\infty}} + \| \lambda^{(k)} \|_{L^{\infty}} \leq 2 \overline{c} \left( |\lambda^0| + 1 + \delta^{1/2} + K + \delta^{-1/2} K\right).
\]

% Then, we find using H\"older's and Young's inequality for a constant $\overline{c}$ (whose value may change from line to line)
% \begin{align}
% &\|\lambda^{(k)}\|_{L^{\infty}} \\ \leq &|\lambda^{(k)}(0)| + \overline{c} \int_{0}^T |N^T \gamma^{(k)} + C \mu^{(k)}| + |\sigma^{(k)}| + |\rho^{(k)}| \, dt  \\ \leq &\overline{c} \left(|\lambda^0| + \| N^T \gamma^{(k)} + C \mu^{(k)}\|_{L^2} +  \left( \int_0^T \frac{|\sigma^{(k)}|^2 + |\rho^{(k)}|^2 }{1+|\lambda^{(k)}| + |\gamma^{(k)}|} \, dt \right)^{1/2} (1 + \| \lambda^{(k)} \|_{L^{\infty}} + \| \gamma^{(k)} \|_{L^{\infty}})^{1/2} \right)\\
% \leq &\overline{c} \left(|\lambda^0| + (K + \delta \| \gamma^{(k)}\|_{L^{\infty}})^{1/2} +  \left( K + \delta \| \gamma^{(k)}\|_{L^{\infty}} \right)^{1/2}  (1 + \| \lambda^{(k)} \|_{L^{\infty}} + \| \gamma^{(k)} \|_{L^{\infty}})^{1/2}\right) \\
% \leq &\overline{c} \left( |\lambda^0| + 1 +  K + \delta \| \gamma^{(k)}\|_{L^{\infty}} + \frac{K}{\eps} + \frac{\delta}{\eps} \| \gamma^{(k)} \|_{L^{\infty}} + \eps ( 1 + \| \lambda^{(k)} \|_{L^{\infty}} + \| \gamma^{(k)} \|_{L^{\infty}}) \right).
% \end{align}
% Next, we estimate similarly using $|\gamma^{(k)}(T)| = | \lambda^{(k)}(T) | \leq \|\lambda^{(k)} \|_{L^{\infty}}$ and the estimate above
% \begin{align}
% &\| \gamma^{(k)} \|_{L^{\infty}} \\ \leq &\overline{c} \left( |\gamma^{(k)}(T)| + |\lambda^{(k)}(T)| + \int_0^T |N^T \gamma^{(k)} + C \mu^{(k)}| + |\sigma^{(k)}| + |\rho^{(k)}|\, ds \right) \\
% \leq &\overline{c} \left( \|\lambda^{(k)}\|_{L^{\infty}} +  \| N^T \gamma^{(k)} + C \mu^{(k)}\|_{L^2} + \left( \int_0^T \frac{|\sigma^{(k)}(s)|^2 + |\rho^{(k)}(s)|^2}{1 + |\lambda^{(k)}| +  |\gamma^{(k)}| } \, ds \right)^{1/2}  (1 + \| \lambda^{(k)} \|_{L^{\infty}} + \| \gamma^{(k)} \|_{L^{\infty}})^{1/2}  \right) \\
% \leq &\overline{c} \left(  \|\lambda^{(k)}\|_{L^{\infty}} + 1 +  K + \delta \| \gamma^{(k)}\|_{L^{\infty}} + \frac{K}{\eps} + \frac{\delta}{\eps} \| \gamma^{(k)} \|_{L^{\infty}} + \eps ( 1 + \| \lambda^{(k)} \|_{L^{\infty}} + \| \gamma^{(k)} \|_{L^{\infty}}) \right) \\
% \leq &\overline{c} \left( |\lambda^0| + 1 +  K + \delta \| \gamma^{(k)}\|_{L^{\infty}} + \frac{K}{\eps} + \frac{\delta}{\eps} \| \gamma^{(k)} \|_{L^{\infty}} + \eps ( 1 + \| \lambda^{(k)} \|_{L^{\infty}} + \| \gamma^{(k)} \|_{L^{\infty}}) \right).
% \end{align}
% Hence, if $4 \overline{c} \delta^{1/2} + 2\overline{c} \delta < \frac12$ then choose $\eps = \delta^{1/2}$ to find
% \begin{align*}
% &\| \gamma^{(k)} \|_{L^{\infty}} + \| \lambda^{(k)} \|_{L^{\infty}} \\ \leq &
% 2\overline{c}\left( |\lambda^0| + 1 +  K + \delta \| \gamma^{(k)}\|_{L^{\infty}} + K \delta^{-1/2}+ \delta^{1/2} \| \gamma^{(k)} \|_{L^{\infty}} + \delta^{1/2} ( 1 + \| \lambda^{(k)} \|_{L^{\infty}} + \| \gamma^{(k)} \|_{L^{\infty}}) \right) \\
% \leq &\frac12 \left( \| \gamma^{(k)} \|_{L^{\infty}} + \| \lambda^{(k)} \|_{L^{\infty}} \right) + 2\overline{c} \left( |\lambda^0| + 1 +  K  + K \delta^{-1/2} + \delta^{1/2}\right).
% \end{align*}
In particular, $\lambda^{(k)}$ and $\gamma^{(k)}$ are uniformly bounded in $L^{\infty}$. By \eqref{eq: est coer dual} this implies that the functions $\rho^{(k)}$ and $\sigma^{(k)}$ are uniformly bounded in $L^2$. Consequently, 
\begin{align*} 
&\sup_k \| \dot{\gamma}^{(k)} \|_{L^2} + \| \dot{\lambda}^{(k)} \|_{L^2}  \\ \leq &\sup_k \overline{c} \left( \| \lambda^{(k)} \|_{L^2} + \| \gamma^{(k)} \|_{L^2} + \| \rho^{(k)} \|_{L^2} + \| \sigma^{(k)} \|_{L^2} +  \| N^T \gamma^{(k)} + C \mu^{(k)}\|_{L^2} \right) < \infty.
\end{align*}
This shows that the sequences $(\gamma^{(k)})_k, (\lambda^{(k)})_k$ are uniformly bounded in $H^1((0,T);\R^n)$. The uniform boundedness of $(\mu^{(k)})_k$ in $L^2(\Omega;\R^m)$ then follows from \eqref{eq: est coer dual}. 

Let us now consider the special case $F=0$. 
First, note that by Proposition \ref{prop: lb g} it holds in the notation from above similarly to \eqref{eq: est coer dual} that
\begin{equation} \label{eq: est dual F0}
\int_0^T \frac1{2a_u} |\sigma^{(k)}|^2 + \frac1{2a_p} |\rho^{(k)} |^2 + \frac1{2a_u} |N^T\gamma^{(k)} + C \mu^{(k)}|^2  \, dt \leq K + (|x^0| + \| A \|_{L^1}) \| \gamma^{(k)} \|_{L^{\infty}}. 
\end{equation}
By \eqref{eq: est gamma0} it follows that
\[
|\gamma^{(k)}(0)| \leq \overline{c} \left( |\lambda^0| + ( K + (|x^0| + \| A \|_{L^1}) \| \gamma^{(k)} \|_{L^{\infty}})^{1/2} \right).
\]
In turn, using \eqref{eq: sol ode} and \eqref{eq: est dual F0} it follows
\begin{align}
     &\| \gamma^{(k)} \|_{L^{\infty}} + \| \lambda^{(k)} \|_{L^{\infty}} \\
    \leq &\overline{c} \left( |\gamma^{(k)}(0)| +  |\lambda^0| + \int_0^T |\sigma^{(k)}| + |\rho^{(k)}| + |N^T\gamma^{(k)} + C \mu^{(k)}| \, ds \right) \\
    \leq &\overline{c} \left( |\lambda^0| + ( K + (|x^0| + \| A \|_{L^1}) \| \gamma^{(k)} \|_{L^{\infty}})^{1/2} \right).
\end{align}
This implies that $\sup_k \| \gamma^{(k)} \|_{L^{\infty}} + \| \lambda^{(k)} \|_{L^{\infty}} < \infty$. The rest of the proof then concludes analogously to before.

% Let us now assume that $F=0$. In this case we obtain from Proposition \ref{prop: lb g} with the notation above similarly to \eqref{eq: est coer dual} that 
% \[
% \| \sigma^{(k)} \|_{L^2}^2 + \| \rho^{(k)} \|_{L^2}^2 + \| N^T \gamma^{(k)} - C \mu^{(k)}\|_{L^2}^2 \leq \overline{c} (K + \| \gamma \|_{L^{\infty}} (\| A\|_{L^1} + |x^0|) ).
% \]
% It follows for $\eps > 0$ similarly to above using Young's inequality
% \begin{align*}
% \| \lambda^{(k)} \|_{L^{\infty}} &\leq  \overline{c} \left( |\lambda^0| + \| N^T \gamma^{(k)} - C \mu^{(k)} \|_{L^2} + \| \sigma^{(k)} \|_{L^2} + \| \rho^{(k)} \|_{L^2} \right) \\
% & \leq \overline{c} \left( |\lambda^0| + K^{1/2}  + \eps \| \gamma^{(k)} \|_{L^{\infty}} + \frac1{\eps} (\| A\|_{L^1} + |x^0|) \right)
% \end{align*}
% and 
% \begin{align*}
%     \| \gamma^{(k)} \|_{L^{\infty}} &\leq  \overline{c} \left( |\gamma^{(k)}(T)| + |\lambda^{(k)}(T)| + \| N^T \gamma^{(k)} - C \mu^{(k)} \|_{L^2} + \| \sigma^{(k)} \|_{L^2} + \| \rho^{(k)} \|_{L^2} \right) \\
% & \leq \overline{c} \left( \|\lambda^{(k)} \|_{L^{\infty}} + K^{1/2}  + \eps \| \gamma^{(k)} \|_{L^{\infty}} + \frac1{\eps} (\| A\|_{L^1} + |x^0|) \right) \\
% &\leq \overline{c} \left( K^{1/2}  + \eps \| \gamma^{(k)} \|_{L^{\infty}} + \frac1{\eps} (\| A\|_{L^1} + |x^0|) \right).
% \end{align*}
% Then choosing $\overline{c} \eps = \frac12$, we obtain that $\| \gamma^{(k)}\|_{L^{\infty}}$ and $\| \lambda^{(k)}\|_{L^{\infty}}$ are uniformly bounded. 
% Using \eqref{eq: rep derivative} it follows that
% \begin{align*}
% \| \dot{\gamma}^{(k)} \|_{L^2} + \| \dot{\lambda}^{(k)} \|_{L^2} &\leq \overline{c} \left( \| \gamma^{(k)} \|_{L^\infty} + \| \lambda^{(k)} \|_{L^{\infty}}  + |\lambda^{(k)}(T)| + \| N^T \gamma^{(k)} - C \mu^{(k)} \|_{L^2} + \| \sigma^{(k)} \|_{L^2} \right) \\ &\leq \overline{c} (K^{1/2} +  \| \gamma^{(k)} \|_{L^{\infty}} + \| \lambda^{(k)} \|_{L^{\infty}}) + \overline{c} \| \gamma^{(k)} \|_{L^{\infty}}^{1/2} (\| A\|_{L^1} + |x^0|)^{1/2} .
% \end{align*}
% Consequently, the sequences $(\gamma^{(k)})_k$ and $(\lambda^{(k)})_k$ are bounded in $H^1$.
\end{proof}

\end{document}
